%-------------------------------------------------------
%
% These are the shadow realms; nothing here is included
%
%-------------------------------------------------------


\section{Problem Setting}

Let $\Theta \subset \R^{p}$ and $\Theta' \subset \R^q$ denote the geometric and physical parameter spaces respectively. Further we let $\Omega \subset \R^d$ be some Lipschitz domain and $\Gamma = [0, T]$ for some $T > 0$ be the trial space of a parametric time-dependent PDE, i.e. $u(t; \mu, \nu) = u(x, t; \mu, \nu)$ is the solution of the PDE, if
\begin{equation*}\begin{cases}
    \partial_t u(t; \mu, \nu) + \mathcal{L}(\mu, \nu)u(t; \mu, \nu) + \mathcal{N}(u(\mu, \nu), \mu, \nu) = f(t; \mu, \nu), \quad &(x, t) \in \Omega \times \Gamma, \\
    \mathcal{B}(\mu, \nu) u(t; \mu, \nu) = g_N(t; \mu, \nu) &(x, t) \in \partial \Omega \times \Gamma, \\
    u(t; \mu, \nu) = u_0(\mu, \nu) &(x, t) \in \Omega \times \{0 \}. \\
\end{cases}
\end{equation*}
Assume we half-discretize the problem in the spacial domain $\Omega$ with $N_h$ degrees of freedom. Note that we refer to this discrete grid space on $\Omega$ as $\Omega^h$ and its boundary grid points as $\partial \Omega^{h}$. Then the resulting full order model can be seen as the solution to 
\begin{equation*} \begin{cases}
    M(\mu, \nu) \partial_{t} u_h(t; \mu, \nu) + A(\mu, \nu) u_h(t; \mu, \nu) + N(u_h(\mu, \nu), \mu, \nu)= f(t; \mu, \nu)  \quad &(x, t) \in \Omega^{h} \times \Gamma, \\
    u_h(0; \mu, \nu) = u_0(\mu, \nu) \quad &(x, t) \in \Omega^{h} \times \{0\}, \end{cases} 
\end{equation*}
where we assume $M: \Theta \times \Theta' \to \R^{N_h \times N_h}$ to send the parameter set to a mass matrix, $A: \Theta \times \Theta' \to \R^{N_h \times N_h}$ to a symmetric positive definite stiffness matrix, $N: \R^{N_h} \to \R^{N_h}$ to be some non-linear term and $f: \Gamma \times \R^{N_h} \times \Theta \times \Theta' \to \R^{N_h}$ to be the corresponding source term of the equation. Additionally we assume to have initial conditions via the function $u_0: \Theta \times \Theta' \to \Omega^{N_h}$.

The goal is to find the parameter-to-solution map $\mathcal{F}: \Theta \times \Theta \to \R^{N_h} \times \Gamma$, such that for all potential pairs $(\mu, \nu)$ the resulting $\hat{u}_{\mu, \nu}(t) \in \R^{N_h}$ is "close" to $u_h(t; \mu, \nu)$ in an $L^2$-sense, i.e. as a projection onto the basis functions of the discretized solution manifold.

\subsection{Example 1: A time-dependent convection-diffusion Problem with inflow}

For our deployment of the DOD-DL-ROM one good example might be the time-dependent variant of the convection-diffusion equation with sharp inflow boundary terms, since the wanted KnW properties were given in that case and its a good bet, that here the evolution in time will most likely only add few dimensions along the $\mu$-invariant trial manifold. Thus we define with $\Theta = [0.2, \pi - 0.2]$ and $\Theta' = [0.5, 1]$ the parametric PDE, which we will discretize and then solve via a full order method,
\begin{equation*} \begin{cases}
    \partial_t u(t; \mu, \nu) = \mbox{div} \left( A(\nu) \nabla u(t; \mu, \nu) \right) + b(\mu) \cdot \nabla u(t; \mu, \nu) \qquad &(x, t) \in \Omega \times \Gamma, \\
    \nabla u(t; \mu, \nu) \cdot \mathbf{n} = u_0, \qquad &(x, t) \in \partial \Omega \times \Gamma, \\
    \end{cases}
\end{equation*}
where $A(\nu) = \nu, \, b(\mu) = 30 \cdot [ \cos \mu, \, \sin \mu ]^T$ and
\begin{equation*}
    u_0(x, y) = \begin{cases} -(y - 0.5)^2 + 0.25, \quad &\text{for } x = 0, \\
    0, \quad &\text{else.} \end{cases}
\end{equation*}

\subsection{Example 2: Another time-dependent convection-diffusion Problem}
We now transfer the interesting and simpler parts of the problem presented in \cite{POD-DL-ROM}. Let $\Theta = [0.3, 0.7 ] \times [0.3, 0.7] \times [30, 70]$ and $\Theta' = [0.002, 0.005]$ be the parameter spaces for the following equation, for which we denote $\mu = [\mu_1, \mu_2, \mu_3]^T$
\begin{equation*} \begin{cases}
    \partial_t u(t; \mu, \nu) -  \mbox{div}\left( A(\nu) \nabla u(t; \mu, \nu) \right) + b(\mu_1) \cdot \nabla u(t; \mu, \nu) = f(\mu_2, \mu_3) \qquad &(x, t) \in \Omega \times \Gamma, \\
    \nabla u(t; \mu, \nu) \cdot \mathbf{n} = 0 \qquad &(x, t) \in \partial \Omega \times \Gamma, \\
    u(t; \mu, \nu) = 0 \qquad &(x, t) \in \Omega \times \{ 0\}, \\
    \end{cases}
\end{equation*}
where we define $A(\nu) = \nu$,
\begin{equation*}
    f(x, y; \mu_2, \mu_3) = 10 \exp \left(-((x - \mu_2)^2 + (y - \mu_3)^2) / 0.07^2\right),
\end{equation*}
and
\begin{equation*}
    b(\mu_1) = [\cos(\pi / \mu_1), \sin(\pi / \mu_1)]^T.
\end{equation*}


\section{Architecture}

Let $N_s = N_{train}^2 N_t$. Assuming we have some parameter sample matrix $M \in \R^{(n_{\mu} + 1) \cdot (n_{\nu} + 1) \times N_t}$ and a snapshot matrix $S \in \R^{N_h \times N_s}$ defined via
\begin{gather*}
    M = \left[(t^1, \mu_1, \nu_1) \vert \dots \vert (t^{N_t}, \mu_1, \nu_1) \vert \dots \vert (t^1, \mu_{N_{train}}, \nu_{N_{train}}) \vert \dots \vert (t^{N_t}, \mu_{N_{train}}, \nu_{N_{train}}) \right], \\
    S = \left[u_h(t^1; \mu_1, \nu_1) \vert \dots \vert u_h(t^{N_t}; \mu_1, \nu_1) \vert \dots \vert u_h(t^1; \mu_{N_{train}}, \nu_{N_{train}}) \vert \dots \vert u_h(t^{N_t}; \mu_{N_{train}}, \nu_{N_{train}}) \right],
\end{gather*}
the goal of the master thesis will be to present the DOD-DL-ROM, as a continuous time dependent extension of the paper \cite{DOD} under consideration of continuous dynamics under parameteric influence, thus additionally employing CoLoRA layers as discussed in \cite{CoLoRA}. The original POD-DL-ROM proposed in \cite{POD-DL-ROM} is the basis of this idea.

\subsection{General DL-ROM}

A general Deep learning based reduced order model (DL-ROM) firstly appeared in architecture and training in the paper \cite{DL-ROM}; to properly approximate the potentially non-linear solution manifold of the high fidelity FOM, we employ an Autoencoder (AE), and on the reduced manifold we employ some other Deep Learning architecture, like the proposed Deep Feed-Forward Neural Network (DFNN), to learn the latent dynamics of the parametric PDE. This concludes in the following approximation of the FOM solution $u_h$
\begin{equation}
    \hat{u}_h(t; \mu, \nu) = \Psi_h(\Phi_n(t; \mu, \nu, \theta_{DF}), \theta_{D}) \in \R^{N_h},
\end{equation}
where we set $\Psi_h = f^D_h : \R^n \times \Theta_D \to \R^{N_h}$ for the decoder (and $f^E_n: \R^{N_h} \times \Theta_E \to \R^n$ the corresponding encoder) of a convolutional AE. Then $f^E_n$ denotes the projection onto the reduced trial solution manifold and we assume $\Phi_n = \Phi_n^{DF} : \Gamma \times \Theta \times \Theta' \times \Theta_{DF} \to \R^n$ to be the DFNN mapping the dynamics of the parametric PDE onto the reduced trial solution manifold.

The subsequent training is then performed via a loss function of a convex combination of both, the projection error using the encoder and the dynamics error considering both the decoder and the feed forward network, i.e. for $\theta = (\theta_D, \theta_E, \theta_{DF})$ define the optimal weights via
\begin{align}
    \theta^* = \argmin_\theta \frac{1}{N_s} \sum_{i, j = 1}^{N_{train}} \sum_{k = 1}^{N_t} \, &\frac{\omega_h}{2} \, \norm{u_h(t^k; \mu_i, \nu_j) - f^D_h(\Phi_n^{DF}(t^k; \mu_i, \nu_j, \theta_{DF}), \theta_D), \theta_{D})}^2 \\
    + \, &\frac{1 - \omega_h}{2} \, \norm{f^E_n(u_h(t^k; \mu_i, \nu_j), \theta_E) - \Phi_n^{DF}(t^k; \mu_i, \nu_j, \theta_{DF})}^2.
\end{align}

\subsection{DOD-DL-ROM}

For comparison, one can see, that the POD-DL-ROM works significantly more efficient than the DL-ROM by pre-reducing the dimension $N_h$ to some $N < N_h$ in a linear POD-based fashion. This has the disadvantage of potentially losing fidelity, if the Kolmogorov $N$-width is rather large, which is the case for non-linear PDEs, relying on geometric circumstances. For this, we might instead of using snapshot-based rPOD for the matrix $S$ (optionally the correlation matrix) rather directly employ a Deep Orthogonal Decomposition to find the reduction matrix in dependency of $\mu \in \Theta$ and $t \in \Gamma$.

Thus we have a common denominator for the DOD-DL-ROM approach. Defining $V: \Gamma \times \Theta \times \Theta_{DOD} \to \R^{N_h \times N}; V(t; \mu, \theta_{DOD}) \mapsto \A \Tilde{V}_{N_A}(t; \mu, \theta_{DOD})$, where $\A \in \R^{N_h \times N_A}$ is the pre-reduction POD for $N < N_A < N_h$ employed to deal with exceptionally large inherent dimensions $N_h$ and $\Tilde{V}_{N_A}(t; \mu, \theta_{DOD}) \in \R^{N_A \times N}$ is called \emph{inner module} of the DOD, yields the foundation of the reduction consideration. From hereon, there are a couple of different approaches, which lead to similar ideas; one is able to, after learning the DOD-map, either suppose, the remaining reduced solution manifold is "linear" enough to work directly with a coefficient finder network (Option A) or a further non-linear reduction, potentially by an Autoencoder, is needed (Option B). As an additional, maybe more ambiguous regarding error estimates, we want to introduce a CoLoRA-based ROM, which first tries to use the DOD-NN reductor for the stationary equivalent of the problem (if existent) and then employs a CoLoRA architecture to reduce the dimension up to the dynamics of the system (Option C). 

For Options A and B: To use the same loss function as in the (POD-)DL-ROM, we do need to guarantee a previous construction of the inner module of the DOD. This first step is done with a separate training with a time dependent analogue to the DOD, i.e.
\begin{equation}
    \theta_{DOD}^* = \argmin_{\theta_{DOD}} \frac{1}{N_s} \sum_{i,j = 1}^{N_{train}} \sum_{k = 1}^{N_t} \, \norm{u_h(t^k; \mu_i, \nu_j) - V(t^k; \mu_i, \theta_{DOD}) V^T(t^k; \mu_i, \theta_{DOD}) \G u_h(t^k; \mu_i, \nu_j)}^2.
\end{equation}

\subsubsection{Linear DOD-DL-ROM (Option A)}

Let $\Phi_n^{DLNN}: \Gamma \times \Theta \times \Theta' \times \Theta_{DLNN} \to \R^{N}$ be some Deep Feed-Forward Neural Network, then the overall workflow can be summed up by
\begin{equation}
    \hat{u}_h(t; \mu, \nu, \theta_{DOD}, \theta_{DFNN}) = \A \Tilde{V}_{N_A}(t; \mu, \theta_{DOD}) \cdot \Phi_n^{DFNN}(t; \mu, \nu, \theta_{DFNN}).
\end{equation}
This is a rather obvious extension of \cite{DOD}, but works surprisingly well. It will be the basic benchmark to overcome for the following algorithms. As for training a general comparison between the projection of a true solution and the prediction of the Feed Forward Network works well, i.e. for $\theta = \theta_{DFNN}$ let the optimal weights be given by
\begin{equation}
    \theta^* = \argmin_\theta \frac{1}{N_s} \sum_{i, j = 1}^{N_{train}} \sum_{k = 1}^{N_t} \norm{V^T(t^k; \mu_i, \theta^*_{DOD}) u_h(t^k; \mu_i, \nu_j) - \Phi_n^{DFNN}(t^k; \mu_i, \nu_j, \theta_{DFNN})}.
\end{equation}

\subsubsection{Non-linear DOD-DL-ROM (Option B)}

Let $f_N^D$ and $f_n^E$ be the decoder and encoder of an AE respectively as above. Each of these non-linear DOD-DL-ROMs is subject to further dimensional reduction and thus is not comparable directly to the linear method above. Nonetheless in terms of reduction in dimensionality both methods can be comparable, if the latent dynamics dimension $n$ is the same in both cases. The general workflow is given by
\begin{equation}
    \hat{u}_h(t; \mu, \nu, \theta_{DOD}, \theta_D, \theta_{DFNN}) = \A \Tilde{V}_{N_A}(t; \mu, \theta_{DOD}) \cdot f_N^D(\Phi_n^{DFNN}(t; \mu, \nu, \theta_{DFNN}), \theta_{D}).
\end{equation}

Since for this, we employ an AE, we optimize for two loss functions, since the employed decoder is tied strictly to its respective encoder being able to correctly project onto the low-dimensional subspace of the latent dynamics of the PDE. For $\theta = (\theta_D, \theta_E, \theta_{DFNN})$ define the optimal weights via
\begin{align*}
    \theta^* = &\argmin_\theta \frac{1}{N_s} \sum_{i, j = 1}^{N_{train}} \sum_{k = 1}^{N_t} \\
    &\frac{\omega_h}{2} \, \norm{u_h(t^k; \mu_i, \nu_j) - V(t^k; \mu_i, \theta_{DOD}^*) \cdot f^D_h(\Phi_n^{DFNN}(t^k; \mu_i, \nu_j, \theta_{DFNN}), \theta_D), \theta_{D})}^2 \\
    + \, &\frac{1 - \omega_h}{2} \, \norm{V^T(t^k; \mu_i, \theta_{DOD}^*) \cdot f^E_n(u_h(t^k; \mu_i, \nu_j), \theta_E) - \Phi_n^{DFNN}(t^k; \mu_i, \nu_j, \theta_{DFNN})}^2.
\end{align*}

\subsection{CoLoRA-DL (Option C)}

Here we aim at minimizing the dimension through two major reductions; the first being the reduction of the stationary DOD to the latent dimension of the $\mu$-invariant solution manifold of the stationary equivalent to the problem (regarding the Dirichlet boundary conditions as the initial conditions on some a priori defined boundary), and the second concerning the actual reduction with CoLoRA-layered architecture to the latent dynamic dimension $n << N_h$. The general workflow is thus given by 
\begin{equation}
    \hat{u}_h(t; \mu, \nu, \theta^*_{sDOD}, \theta_{C}) = \A \cdot \mathcal{C}_L( \dots (\mathcal{C}_2(\mathcal{C}_1(V_0(\mu, \theta^*_{sDOD}) \cdot \phi_0(\mu, \nu, \theta^*_{sDOD})), \theta_{C_1}), \theta_{C_2}), \theta_{C_L}),
\end{equation}
where $\mathcal{C}_i: \R^{N_A} \times \Theta_C \to \R^{N_A}; \mathcal{C}_i(x, \theta_{C_i}) = W_i x + A_i \, \text{diag}(\alpha_i(\nu, t)) B_i x + b_i$ for all $1 \leq i \leq L.$
Remark, that we will still employ a pre-reduction using the rPOD matrix $\A$ stemming from the same snapshot matrix $S$ as in both Options before.





% Requires Hölder definition...

Another interesting result is the famous Embedding Theorem, that may give the reader an intuition for
the relation between Sobolev and Hölder spaces. We will not employ it in this thesis, however it 

\begin{theorem}[Sobolev Embedding Theorem]
    Let $m_1, m_2 \in \N_0$, $1 \leq p < \infty$ and $\alpha \in (0, 1)$. If we choose 
    \begin{equation*}
        m_1 - \frac{d}{p} = m_2 + \alpha,
    \end{equation*}
    then the identity defines a continuous embedding
    \begin{equation*}
        W^{m_1, p}(\Omega) \hookrightarrow C^{m_2, \alpha}(\bar{\Omega}).
    \end{equation*}
    If even $m_1 - \frac{d}{p} > m_2 + \alpha$, then the embedding is compact.
\end{theorem}
\begin{proof}
    See \cite[p. 10-13]{alt2016}.
\end{proof}



% Unnecessary


\begin{lemma}[Projection Error] \label{lemma: optimal projection error}
    Let $N_{s_2}$ be given under Assumption (\ref{assumption sample}), then for all $\varepsilon > 0$ there is a ReLU neural network
    $V: \R^{p+1} \to \R^{N_h \times N}$, such that for $u_j := u(\mu, \nu_j, t)$ with given $\mu \in \Theta$ and $t \in \Tcal$
    \begin{equation*}
        \sup\limits_{(\mu, t) \in \Theta \times \Tcal} \bigg\vert \, \frac{1}{N_{s_2}} \sum_{j=1}^{N_{s_2}} \norm{u_j - \V_{\mu, t} \V_{\mu, t}^T u_j}^2 \\
        - \min\limits_{W \in \R^{N_h \times N}} \frac{1}{N_{s_2}} \sum_{j = 1}^{N_{s_2}} \norm{u_j - W W^T u_j}^2 \, \bigg\vert^{1/2} < \varepsilon.
    \end{equation*}
\end{lemma}
\begin{proof}
    The existence and identity of $V^*$ s.t. 
    \begin{equation*}
        V^* = \underset{W \in \R^{N_h \times N}}{\operatorname{argmin}} \sum_{j=1}^{N_{s_2}}
        \norm{u_j - V^* (V^*)^T u_j}^2    
    \end{equation*}
    is given via the statement in Proposition (\ref{prop: eigenvalues and optimal projection under sampling}) with the POD basis.
    Hence defining a norm via 
    \begin{equation*}
        \norm{(u_j)_{j=1}^{N_{s_2}}}_{train} := \bigg\vert \, \sum_{j=1}^{N_{s_2}} \norm{u_j}^2 \, \bigg\vert^{1/2}
    \end{equation*}
    gives us the following inequality firstly using the monotonicity of the square root and then using the triangle inequalities multiple times
    \begin{align*}
        & \, N_{s_2}^{-1/2} \sup\limits_{(\mu, t) \in \Theta \times \Tcal} \bigg\vert \, \sum_{j=1}^{N_{s_2}} \norm{u_j - \V_{\mu, t} \V_{\mu, t}^T u_j}^2 
        - \sum_{j = 1}^{N_{s_2}} \norm{u_j - V^* (V^*)^T u_j}^2 \, \bigg\vert^{1/2} \\
        \leq& \, N_{s_2}^{-1/2} \sup\limits_{(\mu, t) \in \Theta \times \Tcal} \bigg\vert \, \sum_{j = 1}^{N_{s_2}} \norm{\V_{\mu, t} \V_{\mu, t}^T u_j - V^* (V^*)^T u_j}^2 \, \bigg\vert^{1/2}, 
    \end{align*}
    then define $q_j := (V^*)^T u_j$ and $\hat{q}_j := \V_{\mu, t}^T u_j$ and their training vectors accordingly, i.e. $q := (q_j)_{j=1}^{N_{s_2}}, \hat{q} := (\hat{q}_j)_{j=1}^{N_{s_2}}$
    \begin{align*}
        \leq& \, N_{s_2}^{-1/2} \sup\limits_{(\mu, t) \in \Theta \times \Tcal} \norm{\V_{\mu, t} \hat{q} - \V_{\mu, t} q}_{train} + \norm{\V_{\mu, t} q - V^* q}_{train} \\
        \leq& \, N_{s_2}^{-1/2} \sup\limits_{(\mu, t) \in \Theta \times \Tcal}  \norm{\V_{\mu, t}}_\infty \, \norm{\hat{q} - q}_{train} + \norm{\V_{\mu, t} - V^*}_\infty \, \norm{q}_{train} \\
        = & \, N_{s_2}^{-1/2} \sup\limits_{(\mu, t) \in \Theta \times \Tcal} \norm{\V_{\mu, t}}_\infty \, \norm{\V_{\mu, t}^T u - (V^*)^T u}_{train} + \norm{\V_{\mu, t} - V^*}_\infty \, \norm{(V^*)^T u}_{train} \\
        \leq& \, N_{s_2}^{-1/2}  \sup\limits_{(\mu, t) \in \Theta \times \Tcal} \norm{\V_{\mu, t}}_\infty \, \norm{\V_{\mu, t} - V^*}_\infty \, \norm{u}_{train} + \norm{\V_{\mu, t} - V^*}_\infty \, \norm{V^*}_\infty \, \norm{u}_{train} \\
        \leq& \, N_{s_2}^{-1/2}  \sup\limits_{(\mu, t) \in \Theta \times \Tcal} \norm{u}_{train} \left( \norm{\V_{\mu, t} - V^*}_\infty \, \left( \norm{\V_{\mu, t}}_\infty + 1 \right) \right) \\
        \leq& \, 2 M \sup\limits_{(\mu, t) \in \Theta \times \Tcal} \norm{\V_{\mu, t} - V^*}_\infty,
    \end{align*}
    where we used $\norm{u}_{train} \leq \vert \, \sum_{j=1}^{N_{s_2}} M^2 \, \vert^{1/2} = N_{s_2}^{1/2} M$ and additionally $\norm{\V_{\mu, t}}_\infty = \norm{V^*}_\infty = 1$
    by orthonormality in the last arguments.

    \notdone{Do the Yarotski argument with the appropriate number of layers and weights.}

\end{proof}


% copy for safety


\begin{theorem} \label{theo: PAC bound for POD-DL-ROM}
    Let $\mathcal{G}: (\mu, \nu, t) \mapsto \uht$ for any $(\mu, \nu, t) \in \Theta \times \Theta' \times \Tcal$ be the
    parameter-to-solution map and let Assumption (\ref{assumption sample}) and (\ref{assumption parameter-to-solution map}) be true. Let $1  > \delta > 0$ and $\varepsilon > 0$.
    Now assume we can collect freely a sufficient amount of training data $N_{\text{data}} = N_{\text{data}}(\delta, \varepsilon)$ and put it into the
    correlation matrix $K$ as given in Definition (\ref{def: POD}) with $\sigma_k^2$ being its sorted eigenvalues. Now choose
    \begin{equation*}
        N = \operatorname{argmin} \left\{ j \in \N \, \bigg\vert \, \sum_{k > j} \sigma^2_k \leq \frac{m^2}{9} \varepsilon^2 \right\}.
    \end{equation*}
    Given the POD matrix $V \in \R^{N_h \times N}$ we now define the parameter-to-POD-coefficient map 
    $\mathcal{Q}: (\mu, \nu, t) \mapsto q(\mu, \nu, t) :=  V^T \uht$ for 
    any $(\mu, \nu, t) \in \Theta \times \Theta' \times \Tcal$. We assume that there exists $n > 0, \psi_* : \R^n \to \R^N, \psi_*' : \R^N \to \R^n$ that are respectively $s$-times and
    $s'$-times differentiable ( with $s >> s' \geq 2$), such that they enjoy the perfect embedding Assumption (\ref{assumption perfect embedding}), i.e. 
    \begin{equation*}
        \psi_*(\psi_*'(q(\mu, \nu, t))) = q(\mu, \nu, t), \quad \forall (\mu, \nu, t) \in \Theta \times \Theta' \times \Tcal.
    \end{equation*}
    We further let 
    \begin{equation*}
        C_1 := \sup\limits_{\vert \alpha \vert \leq s'} \sup\limits_{v \in \R^N} \vert \, D^\alpha \psi_*'(v) \, \vert, \quad 
        C_2 := \sup\limits_{\vert \alpha \vert \leq s} \sup\limits_{w \in \R^n} \vert \, D^\alpha \psi_*'(w) \, \vert.
    \end{equation*}
    Then there is a constant $c = c(\Theta, \Theta', \Tcal, L, C_1, C_2, p, q, n, s, s')$ and a POD-DL-ROM architecture $V \hat{q} =
    V \psi \circ \phi: \R^{p+q+1} \to \R^{N_h}$ composed of a decoder $\psi: \R^n \to \R^N$ having at most:
    \begin{enumerate}
        \item[-] $L_{n \to N} = c \log(\varepsilon^{-1})$ layers,
        \item[-] $\omega_{n \to N} = cN \varepsilon^{-n/(s-1)} \log(\varepsilon^{-1})$ active weights,
    \end{enumerate}
    and a reduced feed forward neural network $\phi: \R^{p+q+1} \to \R^n$ having at most:
    \begin{enumerate}
        \item[-] $L_{(p+q+1) \to n} = c \log(\varepsilon^{-1})$ layers,
        \item[-] $\omega_{(p+q+1) \to n} = cn \varepsilon^{-(p+q+1)} \log(\varepsilon^{-1})$ active weights,
    \end{enumerate}
    such that $\prob[\Ecal_R < \varepsilon] > 1 - \delta$.
\end{theorem}


% existence and uniqueness subsection for hille yosida

\subsection{Existence and Uniqueness}

This subsection provides a brief introduction to one specific statement of existence and uniqueness, which we coincidentally
can apply well to the projected version of our finite dimensional problem. Note, that of course, this does not justify any existence
result for highly non-linear equations on a lot of upstream classical or even weak problems, like the famous 
incompressible Navier-Stokes Equation.

We start by defining the differential operator class, to which we can then apply the Hille-Yosida Theorem. Most notable, 
we know with that subclass, that any monotone operator, which acts on a finite dimensional space, is in that class, as
will be shown afterwards. 

\begin{definition}[Maximal Monotone Operator] \label{def: maximal monotone operator}
    Let $A: D(A) \subset H \to H$ be an unbounded linear operator. It is \emph{monotone}, if
    \begin{equation*}
        \langle Av, v \rangle_H \geq 0, \quad \text{for all } v \in D(A),
    \end{equation*}
    and \emph{maximal monotone}, if additionally for all $f \in H$ there is $u \in D(A)$ such that
    \begin{equation*}
        u + Au = f.
    \end{equation*}
\end{definition}


\begin{lemma} \label{lemma: finite dimensional monotone operator}
    Let $H$ be a finite dimensional real Hilbert space, and $A: H \to H$ a monotone operator according to Definition
    (\ref{def: maximal monotone operator}). Then $A$ is maximal monotone.
\end{lemma}
\begin{proof}
    By obvious identification with a basis, we know, that $H \cong \R^N$ for some $N \in \N$. Thus any element 
    $u \in H$ can be expressed as 
    \begin{equation*}
        u = \sum_{j=1}^N u_j \varphi_j,
    \end{equation*}
    where $u_j \in \R$ and $\varphi_j \in H$ for all $j=1, \dots, N$, and
    similarly for $i,j, k = 1, \dots, N$ there are $a_{ik} \in \R$ such that
    \begin{equation*}
        Au = \sum_{i=1}^N \sum_{j=1}^N \sum_{k=1}^N a_{ik} u_j \langle \varphi_j, \varphi_i \rangle \varphi_k.
    \end{equation*}
    Thus the condition for $A$ being maximal monotone
    simplyfies to finding $[u_j]_{j=1}^N \in R^N$ for an arbitrary $[f_j]_{j=1}^N \in \R^N$ where
    \notdone{doesnt work}
    \begin{equation*}
        u_j + \sum_{i=1}^N a_{ij} u_j = f_j
    \end{equation*} 
    for all $j=1, \dots, N$, which can obviously be done in general.
\end{proof}


\begin{theorem}[Hille-Yosida Theorem] \label{theo: hille-yosida}
    Let $A$ be a maximal monotone operator, $H$ some Hilbert space and $D(A) \subset H$. 
    Then, given any $u_0 \in D(A)$ there is a unique function
    \begin{equation*}
        u \in C^1([0, \infty); H) \cap C([0, \infty); D(A))
    \end{equation*}
    satisfying
    \begin{equation*}
        \begin{cases}
            \displaystyle
            \frac{\partial u}{\partial t} + Au = 0 \quad &\text{on }[0, \infty), \\
            \displaystyle
            u(0) = u_0 \quad &\text{on }\{0\}.
        \end{cases}
    \end{equation*}
    Moreover,
    \begin{equation*}
        \norm{u(t)}_H \leq \norm{u_0}_H,
    \end{equation*}
    for all $t \geq 0$.
\end{theorem}
\begin{proof}
    See \cite[p. 185-191]{brezis2011}.
\end{proof}



% equivalent estimate to order of N and N', but this time starting from the assumed N = N'

However, we want to start with the focus on the \emph{potential} upper relative error bound. We
remind the reader of the results in Theorem (\ref{theo: relative error POD+DNN}) 
and (\ref{theo: relative error DOD+DNN}). These specify, that if ubiquitously 
$N_s, N_t \to \infty$, $N_A = N_h$ and the complexity of the coefficient searching neural networks
becomes very large, we have
\begin{equation*}
    \Ecal_\text{R}^\text{POD+DNN} \approx \Ecal_{\text{POD}, \infty} = 
    \Ocal\left( N^{-2\alpha} \right) 
    \quad \text{and} \quad
    \Ecal_\text{R}^\text{DOD+DNN} \approx \Ecal_{\text{DOD}, \infty} = 
    \Ocal\left( (N')^{-2\beta} \right).
\end{equation*}
This is of course equivalent to the previous mentioned order of $N$ and $N'$ under the same assumptions.


% vaulted CoLoRA
\subsection{CoLoRA-DL-ROM}

\almostdone{Introduce CoLoRA DL-ROM}

Here we aim at minimizing the dimension through two major reductions; 
the first being the reduction of the stationary DOD to the latent dimension 
of the $\mu$-invariant solution manifold of the stationary equivalent to the problem 
(regarding the Dirichlet boundary conditions as the initial conditions on some a priori defined boundary), 
and the second concerning the actual reduction with CoLoRA-layered architecture to the 
latent dynamic dimension $n << N_h$. The general workflow is thus given by 
\begin{equation}
    \hat{u}_h(t; \mu, \nu, \theta^*_{sDOD}, \theta_{C}) = \A \cdot 
    \mathcal{C}_L( \dots (\mathcal{C}_2(\mathcal{C}_1(V_0(\mu, \theta^*_{sDOD}) 
    \cdot \phi_0(\mu, \nu, \theta^*_{sDOD})), \theta_{C_1}), \theta_{C_2}), \theta_{C_L}),
\end{equation}
where $\mathcal{C}_i: \R^{N_A} \times \Theta_C \to \R^{N_A}; 
\mathcal{C}_i(x, \theta_{C_i}) = W_i x + A_i \, \text{diag}(\alpha_i(\nu, t)) B_i x + b_i$ 
for all $1 \leq i \leq L.$
Remark, that we will still employ a pre-reduction using the rPOD matrix $\A$ stemming from the 
same snapshot matrix $S$ as in both Options before.

% vaulted proof start of concatenations
\begin{proof}
    Intuitively speaking the latter network must account for all "missing" links, which a direct concatenation
    would necessarily be missing; any layer after the "switch" from the first to the latter needs to have a correction
    term in the weights, additionally attributing all potential connections which the first network would impose.
    

    Let w.l.o.g. $L_2 = 2$, since the pass through the first layers of the combined network, is equivalent
    to the pass through the first layers of the individual network. Then we can show by induction:
    \begin{enumerate}
        \item \underline{Base case:} $L_2 = 1$. We argue with $A_{L_1} x = A_{L_1, x_0} \, x$ for $x$ of output dimension of $\Phi^1$, since
        the input dimension of $\Phi^1$ is the output of $\Phi^2$.
        \begin{align*}
            R_\rho(\Phi^1) \circ R_\rho(\Phi^2) (x_0) &= A_{L_1} 
            \left(B_{L_2, x_0} x_0 + B_{L_2, x_1} x_1 + c_{L_2}\right) + b_{L_1} \\
            &= A_{L_1, x_0} B_{L_2, x_0} x_0 + A_{L_1, x_0} B_{L_2, x_1} x_1 + \left(A_{L_1, x_0} c_{L_2} + b_{L_1}\right) \\
            &= A_{L_1, x_0} B_{L_2} \begin{bmatrix}[c|c]x_0^T & x_{1}^T\end{bmatrix}^T + \left(A_{L_1, x_0} c_{L_2} + b_{L_1}\right) \\
            &= R_\rho(\Phi^1 \bullet \Phi^2) (x_0)
        \end{align*}
        \item \underline{Induction assumption:} Let the statement be true for $L_2 \geq 2$ fixed.
        \item \underline{Induction step:} $L_2 \mapsto L_2 + 1$.
    \end{enumerate}
\end{proof}